\section{Consensus Validation}
\label{sec:consensus}

Consensus validation determines whether a node may accept a chain history as
valid. It instantiates the transition relations from
Section~\ref{sec:ledger-model} for headers, block data, active-chain
connection, and reorganization.

\subsection{Validation relation}

Let \(N\) be the network parameter set. For candidate block \(b\), candidate
height \(h\), candidate block time \(t\), active chain \(C\), and UTXO set
\(U\), let \(R=\operatorname{Rules}(N,C,h,t)\). A successful connection
produces:

\[
  (C, U) \xrightarrow{b} (C', U')
\]

where \(C'\) extends or reorganizes \(C\) to include \(b\), and \(U'\) is the UTXO set
obtained after applying every transaction in the connected branch. If any
consensus check fails, the active pair \((C,U)\) is unchanged.

\subsection{Header acceptance}

A header is accepted into the header graph only if all conditions below hold.
Header acceptance does not imply that block transaction data is known or valid.

\begin{longtable}{@{}L{0.28\linewidth}L{0.60\linewidth}@{}}
\toprule
Condition & Required rule \\
\midrule
Uniqueness & If \(\mathrm{bh}(h)\) is already known, the existing header state is reused. \\
Previous header & Every non-genesis header references a known previous header. \\
Proof of work & \(\mathrm{le}(\mathrm{bh}(h)) \le \mathrm{C}^{-1}(h.\text{\code{bits}})\). \\
Target & \(\mathrm{C}^{-1}(h.\text{\code{bits}})\) is valid for the height, network, and difficulty-adjustment interval. \\
Timestamp & The header time is greater than the median timestamp of the previous block and up to ten ancestors. A header whose time is more than \code{MAX\_FUTURE\_BLOCK\_TIME} after the node's current clock time is future-dated and not eligible for active validation until time advances. \\
BIP94 timewarp & On networks enforcing BIP94, a difficulty-adjustment block other than genesis has time at least \code{previous.time - BIP94\_TIMEWARP\_GRACE}. \\
Rule context & \(R\) permits the candidate height, time, version field, and network-specific activation state. \\
Invalid ancestry & A header descending from a known-invalid header is not eligible for best-chain selection. \\
\bottomrule
\end{longtable}

For each accepted header, the node records height, predecessor,
\(\mathrm{C}^{-1}(h.\text{\code{bits}})\), \(\mathrm{bh}(h)\), and accumulated
work. Accumulated work is the sum of the work represented by each header from
genesis through that header.

\subsection{Block-data acceptance}

Full block data is accepted only after the corresponding header is accepted.
The following checks are performed before block data can become a candidate for
active-chain connection:

\begin{longtable}{@{}L{0.26\linewidth}L{0.62\linewidth}@{}}
\toprule
Check & Required rule \\
\midrule
Block form & The block has at least one transaction, satisfies the stripped-size, transaction-count, and weight limits in Section~\ref{sec:blocks}, and the first transaction is the only coinbase transaction. \\
Merkle root & Header \code{merkleRoot} equals \(\mathrm{MR}(L)\) for the block's txid leaves, and the mutated-root condition is rejected. \\
Transaction form & Every transaction passes context-free transaction checks. \\
Legacy sigops limit & \code{legacySigOps(block) * WITNESS\_SCALE\_FACTOR <= MAX\_BLOCK\_SIGOPS\_COST} before P2SH and witness sigops are counted during connection. \\
Contextual finality & Every transaction satisfies \(\operatorname{AbsFinal}\) for the candidate block height and locktime cutoff. \\
Coinbase height & If \(R\) includes BIP34, the coinbase script begins with the candidate height. \\
Witness commitment & If \(R\) includes SegWit and the block has witness data, the coinbase witness commitment is present and correct. \\
Weight & The block weight is at most \code{MAX\_BLOCK\_WEIGHT}. \\
\bottomrule
\end{longtable}

\subsection{Best-chain selection}

Among branches whose headers are valid and whose required block data is
available, the selected branch is the valid branch with greatest accumulated
work. A candidate branch is not eligible while any block on the path from the
active fork to the candidate tip is missing block data or is known invalid.

\begin{enumerate}[leftmargin=*]
  \item Choose the highest-work eligible candidate tip.
  \item Find the fork point with the current active chain.
  \item Disconnect active-chain blocks back to the fork.
  \item Connect candidate-branch blocks forward in height order.
  \item If a candidate block is consensus-invalid, exclude that block and its
        descendants from future selection and retry with the next candidate.
\end{enumerate}

\subsection{Block connection}

Connecting a non-genesis block applies the block's transactions to a temporary
UTXO view whose best block is the previous block hash. The view becomes active
only after all checks below pass.

\begin{longtable}{@{}L{0.28\linewidth}L{0.60\linewidth}@{}}
\toprule
Rule & Required behavior \\
\midrule
BIP30 no-overwrite & A transaction output must not overwrite an existing unspent outpoint except for the historical BIP30 exceptions and buried rules that make the check unnecessary after BIP34. \\
Input availability & Every non-coinbase input references an unspent coin in the candidate view. \\
Duplicate spends & A transaction must not spend the same outpoint more than once. \\
Coinbase maturity & Coinbase-created outputs may be spent only after \code{COINBASE\_MATURITY} confirmations. \\
Sequence locks & Every non-coinbase transaction satisfies the BIP68 relative-lock predicate from Section~\ref{sec:transactions}. \\
Value range & Every output and every sum used during validation remains in the money range. \\
Value conservation & For each non-coinbase transaction, input value is at least output value; the difference is the transaction fee. \\
Script validity & Each input satisfies the previous output's locking script under the script flags active for the block. \\
Sigops cost & Total block signature-operation cost is at most \code{MAX\_BLOCK\_SIGOPS\_COST}. \\
Subsidy & Coinbase output value is at most block subsidy plus total fees. \\
\bottomrule
\end{longtable}

The genesis block is special: its coinbase transaction is part of the block
history but is not connected as a spendable coin by normal block-connection
rules.

\subsection{UTXO transition}

For a valid non-coinbase transaction $tx$, let $I(tx)$ be its input outpoints
and $O(tx)$ be its outputs indexed by output number. The candidate UTXO set is
updated as follows:

\[
  U_{after} =
    (U_{before} \setminus I(tx)) \cup
    \{((txid(tx), n), coin(out_n, height, false)) : out_n \in O(tx)\}.
\]

For the coinbase transaction, there are no spent inputs, and created outputs
are inserted with \code{coinbase = true}. Outputs whose value is zero are still
outputs and are inserted unless another consensus rule rejects the transaction.

\subsection{Rule context}

\(\operatorname{Rules}(N,C,h,t)\) determines the script flags and contextual
block predicates used for connection. Mainnet consensus rule context includes
the buried and deployed rules for P2SH, strict DER signatures, locktime and
sequence verification, SegWit, Taproot, and Tapscript when included in \(R\)
\cite{bip-0016,bip-0065,bip-0066,bip-0068,bip-0112,bip-0113,bip-0141,bip-0143,bip-0147,bip-0341,bip-0342}.

Policy-only flags may be stricter for mempool relay, but policy failure is not
block invalidity.

\subsection{Subsidy and fees}

The block subsidy is determined by height and network parameters. For mainnet,
the subsidy starts at 50 BTC and halves every 210,000 blocks until it reaches
zero by integer right shift \cite{bip-0042}. The consensus rule for a connected
block is:

\[
  \sum outputs(coinbase) \le subsidy(height) + \sum fees(noncoinbase).
\]

Fees are computed from connected inputs and outputs. A transaction whose output
sum exceeds its input sum is invalid.

\subsection{Undo}

Before applying a non-coinbase transaction, a node records the spent coin for
each input in input order. The undo record for a block is the ordered sequence
of those per-transaction records, excluding the coinbase. Undo is sufficient to
restore the previous UTXO set when the block is disconnected; the byte encoding
of undo data is not a consensus rule.

\subsection{Reorganization correctness}

A reorganization from old tip $o$ to new tip $n$ is valid only if:

\begin{enumerate}[leftmargin=*]
  \item the branch ending at $n$ has greater accumulated work than the current
        active branch or is otherwise selected by the same greatest-work rule;
  \item every disconnected block is reversed exactly using the UTXO state that
        existed after that block was connected;
  \item every connected block on the new branch passes normal block connection;
        and
  \item after the transition, the active UTXO set equals the result of applying
        the active chain from genesis to $n$.
\end{enumerate}

Transactions removed from disconnected blocks reenter relay only under current
mempool policy.
